% =============================================================================
% APPENDIX CTM: LIT-AUTHOR: Unknown Researcher Research Integration
% =============================================================================
% Category: LIT
% Language: English
% Version: 1.0 (January 2026)
% Status: Draft
%
% This appendix integrates research papers by Unknown Researcher into the EBF framework.
%
% =============================================================================

\chapter{LIT-AUTHOR: Unknown Researcher Research Integration}
\label{app:ctm}

\begin{abstract}
This appendix integrates the research contributions of Unknown Researcher into the
Evidence-Based Framework. It documents key papers, core findings, and their
integration with EBF dimensions including complementarity parameters ($\gamma$),
context factors ($\Psi$), and utility dimensions.
\end{abstract}

\begin{tcolorbox}[colback=blue!5!white, colframe=blue!75!black,
    title=Quick Reference: Unknown Researcher Research]

\textbf{Appendix:} CTM (LIT)

\textbf{Researcher:} Unknown Researcher

\textbf{Core Themes:}
\begin{itemize}[nosep]
\item Behavioral economics foundations
\item Empirical evidence for EBF parameters
\item Policy implications and interventions
\end{itemize}

\textbf{EBF Integration:}
\begin{itemize}[nosep]
\item Complementarity values ($\gamma$) from experimental evidence
\item Context effects ($\Psi$) documented across studies
\item Utility dimension weightings calibrated to findings
\end{itemize}

\textbf{Cross-References:} BBB (CORE-WHERE), K (LIT-FEHR), G (Glossary)
\end{tcolorbox}

% -----------------------------------------------------------------------------
% APPENDIX SCOPE BOX
% -----------------------------------------------------------------------------

\begin{tcolorbox}[
    colback=orange!5!white,
    colframe=orange!75!black,
    title={\textbf{Appendix Scope: Ziel / In-Scope / Out-of-Scope / Constraints}},
    fonttitle=\bfseries
]

\textbf{Ziel (Objective):}
\begin{itemize}[nosep]
    \item Integrate Unknown Researcher's research into EBF with parameter extraction
\end{itemize}

\textbf{In-Scope:}
\begin{itemize}[nosep]
    \item Paper-by-paper integration with 10C mapping
    \item Extraction of $\gamma$, $\Psi$, and effect size parameters
    \item Cross-references to related EBF components
\end{itemize}

\textbf{Out-of-Scope (delegiert an andere Appendices):}
\begin{itemize}[nosep]
    \item Parameter validation $\rightarrow$ Appendix UNMAPPED_BBB (CORE-WHERE)
    \item Formal axiomatization $\rightarrow$ CORE appendices
    \item Meta-analysis $\rightarrow$ LIT-M appendices
\end{itemize}

\textbf{Constraints (Anwendungsgrenzen):}
\begin{itemize}[nosep]
    \item Parameters are extracted from published studies
    \item Effect sizes may vary across replications
    \item Context-specificity of findings applies
\end{itemize}

\textbf{Lieferobjekte (Deliverables):}
\begin{itemize}[nosep]
    \item Paper integration table with 10C mappings
    \item Extracted parameters for BBB repository
    \item Cross-reference map to related literature
\end{itemize}

\end{tcolorbox}

%%%%%%%%%%%%%%%%%%%%%%%%%%%%%%%%%%%%%%%%%%%%%%%%%%%%%%%%%%%%%%%%%%%%%%%%%%%%%%%
\section{The Fundamental Question}
\label{sec:ctm-fundamental}
%%%%%%%%%%%%%%%%%%%%%%%%%%%%%%%%%%%%%%%%%%%%%%%%%%%%%%%%%%%%%%%%%%%%%%%%%%%%%%%

\begin{tcolorbox}[
    colback=red!5!white,
    colframe=red!75!black,
    title={\textbf{Fundamental Question}}
]
\textbf{How does lit-author: unknown researcher research integration integrate with and contribute to the
Evidence-Based Framework for understanding behavioral outcomes?}

This appendix addresses the core challenge of formalizing behavioral principles
within a rigorous theoretical framework while maintaining practical applicability.
\end{tcolorbox}




%%%%%%%%%%%%%%%%%%%%%%%%%%%%%%%%%%%%%%%%%%%%%%%%%%%%%%%%%%%%%%%%%%%%%%%%%%%%%%%
% APPENDIX AH: Temporal Context (Ψ_T)
% Time Preferences, Discounting, and Intertemporal Choice
%%%%%%%%%%%%%%%%%%%%%%%%%%%%%%%%%%%%%%%%%%%%%%%%%%%%%%%%%%%%%%%%%%%%%%%%%%%%%%%
% Category: CONTEXT
% Language: English

%%%%%%%%%%%%%%%%%%%%%%%%%%%%%%%%%%%%%%%%%%%%%%%%%%%%%%%%%%%%%%%%%%%%%%%%%%%%%%%
\section{Critical Foundations and Objections}
\label{sec:ctm-foundations}
%%%%%%%%%%%%%%%%%%%%%%%%%%%%%%%%%%%%%%%%%%%%%%%%%%%%%%%%%%%%%%%%%%%%%%%%%%%%%%%

\subsection{Objection 1: Scope Limitations}

\textbf{Objection:} The framework presented may have limited applicability
outside the specific domain contexts discussed.

\textbf{Response:} We acknowledge domain-specificity. The framework provides
general principles that should be calibrated to specific contexts. Cross-domain
validation remains an open research question.

\subsection{Objection 2: Measurement Challenges}

\textbf{Objection:} Key constructs may be difficult to measure reliably in
practice.

\textbf{Response:} We recommend using validated scales and instruments where
available, and developing domain-specific proxies where necessary. Sensitivity
analysis should assess robustness to measurement uncertainty.



%%%%%%%%%%%%%%%%%%%%%%%%%%%%%%%%%%%%%%%%%%%%%%%%%%%%%%%%%%%%%%%%%%%%%%%%%%%%%%%
\section{Glossary of Symbols}
\label{sec:ctm-glossary}
%%%%%%%%%%%%%%%%%%%%%%%%%%%%%%%%%%%%%%%%%%%%%%%%%%%%%%%%%%%%%%%%%%%%%%%%%%%%%%%

For comprehensive definitions, see Appendix G (Master Glossary).

\begin{table}[htbp]
\centering
\caption{Key Symbols in Appendix UNMAPPED_CTM}
\begin{tabular}{lll}
\toprule
\textbf{Symbol} & \textbf{Meaning} & \textbf{Reference} \\
\midrule
$\gamma$ & Complementarity parameter & Appendix UNMAPPED_B \\
$\Psi$ & Context vector & Appendix UNMAPPED_V \\
$U$ & Utility function & Chapter 3 \\
\bottomrule
\end{tabular}
\end{table}



%%%%%%%%%%%%%%%%%%%%%%%%%%%%%%%%%%%%%%%%%%%%%%%%%%%%%%%%%%%%%%%%%%%%%%%%%%%%%%%
\section{Open Issues and Research Agenda}
\label{sec:ctm-open}
%%%%%%%%%%%%%%%%%%%%%%%%%%%%%%%%%%%%%%%%%%%%%%%%%%%%%%%%%%%%%%%%%%%%%%%%%%%%%%%

\begin{enumerate}
    \item \textbf{Empirical Validation:} Further testing across diverse contexts
    \item \textbf{Parameter Refinement:} Improved estimation methods needed
    \item \textbf{Integration:} Enhanced cross-appendix linkages
\end{enumerate}

\section{Temporal Context: Time Preferences and Intertemporal Choice}

\nocite{bcm_master}

\label{app:temporal-context}

%==============================================================================
% PART 0: INTEGRATION OVERVIEW
%==============================================================================
\subsection{Framework Integration Map}

This appendix fills a critical gap in the $(C, \Psi, C^*, K, Q)$ framework: 
the temporal dimension $\Psi_T$. How people trade off present and future 
profoundly shapes economic behavior---savings, investment, health, education, 
addiction, and environmental policy all depend on time preferences. The 
behavioral revolution revealed that human discounting is not exponential 
(as standard theory assumes) but hyperbolic, leading to time inconsistency, 
present bias, and self-control problems. This has fundamental implications 
for $C^*$ and welfare $Q$.

\begin{center}
\renewcommand{\arraystretch}{1.4}
\begin{tabular}{|l|l|l|}
\hline
\textbf{Framework Element} & \textbf{Temporal Contribution} & \textbf{Key Papers} \\
\hline
$\Psi_T$ (time preferences) & Discount rates, patience & Frederick et al. 2002 \\
Present bias & $\beta < 1$ immediate gratification & Laibson 1997, O'Donoghue-Rabin 1999 \\
Time inconsistency & Plans change over time & Strotz 1956, Thaler-Shefrin 1981 \\
Commitment devices & $\Psi_{design}$ for self-control & Ashraf et al. 2006, Bryan et al. 2010 \\
Long-run decisions & Savings, health, education & Madrian-Shea 2001, Chetty et al. 2014 \\
Intergenerational & Climate, debt, institutions & Stern 2007, Nordhaus 2007 \\
\hline
\end{tabular}
\end{center}

\paragraph{Why $\Psi_T$ Matters.}
Temporal context shapes nearly every economic decision:
\begin{enumerate}
    \item \textbf{Savings}: Present bias $\Rightarrow$ undersaving for retirement
    \item \textbf{Health}: Immediate pleasure vs. long-term health
    \item \textbf{Education}: Costly now, benefits later
    \item \textbf{Addiction}: Hyperbolic discounting explains relapse
    \item \textbf{Environment}: Discounting determines climate policy
    \item \textbf{Policy}: Commitment devices can improve $Q$
\end{enumerate}

%==============================================================================
% PART I: FOUNDATIONS – EXPONENTIAL DISCOUNTING
%==============================================================================
\subsection{Part I: Foundations – The Standard Model}

\subsubsection{Samuelson's Discounted Utility}

Samuelson (1937) introduced the standard model:

\begin{equation}
\boxed{
U = \sum_{t=0}^{T} \delta^t u(c_t) \quad \text{where } \delta = \frac{1}{1+\rho}
}
\label{eq:exponential}
\end{equation}

\paragraph{Properties of Exponential Discounting.}
\begin{itemize}
    \item \textbf{Time consistency}: Plans made today are followed tomorrow
    \item \textbf{Constant discount rate}: $\rho$ same for all horizons
    \item \textbf{Stationarity}: Only time difference matters, not calendar time
\end{itemize}

\subsubsection{Framework Translation}

Standard model assumes $\Psi_T = \delta$ constant:
\begin{equation}
\Psi_T^{standard} = \delta \quad \text{(one parameter)}
\end{equation}

\subsubsection{Problems with the Standard Model}

Frederick, Loewenstein \& O'Donoghue (2002) documented anomalies:
\begin{itemize}
    \item Discount rates vary by domain (money vs. health)
    \item Discount rates decline with horizon (hyperbolic)
    \item Sign effect: Losses discounted less than gains
    \item Magnitude effect: Small amounts discounted more
    \item Delay-speedup asymmetry
\end{itemize}

%==============================================================================
% PART II: HYPERBOLIC DISCOUNTING AND PRESENT BIAS
%==============================================================================
\subsection{Part II: Hyperbolic Discounting – The Behavioral Revolution}

\subsubsection{Strotz: Time Inconsistency}

Strotz (1956) first identified time inconsistency:

\begin{equation}
\boxed{
\text{Optimal plan at } t \neq \text{Optimal plan at } t' > t
}
\label{eq:strotz}
\end{equation}

With non-exponential discounting, preferences reverse over time.

\paragraph{Example.}
\begin{itemize}
    \item Today: Prefer \$110 in 31 days over \$100 in 30 days
    \item In 30 days: Prefer \$100 today over \$110 tomorrow
\end{itemize}

Same tradeoff (1 day, \$10), different choices!

\subsubsection{Hyperbolic Discounting}

Ainslie (1975) and Loewenstein \& Prelec (1992) formalized:

\begin{equation}
\boxed{
D(t) = \frac{1}{1 + kt} \quad \text{(hyperbolic)}
}
\label{eq:hyperbolic}
\end{equation}

Compare to exponential: $D(t) = \delta^t$

\paragraph{Key Difference.}
Hyperbolic: Discount rate $\frac{D'(t)}{D(t)} = \frac{k}{1+kt}$ \textit{declines} with $t$.

Exponential: Discount rate $-\ln(\delta)$ \textit{constant}.

\subsubsection{Framework Translation}

Hyperbolic discounting means $\Psi_T$ is time-varying:
\begin{equation}
\Psi_T(t) = \frac{1}{1+kt} \quad \text{(depends on delay)}
\end{equation}

\subsubsection{Laibson: Quasi-Hyperbolic ($\beta$-$\delta$) Model}

Laibson (1997) introduced tractable approximation:

\begin{equation}
\boxed{
U = u(c_0) + \beta \sum_{t=1}^{T} \delta^t u(c_t) \quad \text{where } \beta < 1
}
\label{eq:beta-delta}
\end{equation}

\paragraph{Interpretation.}
\begin{itemize}
    \item $\beta$: Present bias (immediate gratification premium)
    \item $\delta$: Long-run patience (standard discounting)
    \item $\beta = 1$: Standard exponential model
    \item $\beta < 1$: Extra weight on present
\end{itemize}

\paragraph{Empirical Estimates.}
\begin{itemize}
    \item $\beta \approx 0.7$ (substantial present bias)
    \item $\delta \approx 0.99$ (modest long-run impatience)
\end{itemize}

\subsubsection{Framework Translation}

Present bias is captured by:
\begin{equation}
\Psi_T = (\beta, \delta) \quad \text{where } \beta < 1 \text{ captures present bias}
\end{equation}

%==============================================================================
% PART III: SOPHISTICATION VS. NAIVETE
%==============================================================================
\subsection{Part III: Sophistication vs. Naïveté}

\subsubsection{O'Donoghue \& Rabin: Self-Awareness}

O'Donoghue \& Rabin (1999, 2001) distinguished:

\begin{equation}
\boxed{
\text{Sophisticates: Know } \beta < 1 \quad | \quad \text{Naïfs: Believe } \beta = 1
}
\label{eq:sophistication}
\end{equation}

\paragraph{Sophisticates.}
\begin{itemize}
    \item Predict future self-control failures
    \item May seek commitment devices
    \item Can still procrastinate (knowing they will)
\end{itemize}

\paragraph{Naïfs.}
\begin{itemize}
    \item Believe they'll behave patiently tomorrow
    \item Perpetually postpone costly actions
    \item Repeatedly surprised by own behavior
\end{itemize}

\subsubsection{Partial Naïveté}

O'Donoghue \& Rabin (2001) introduced:
\begin{equation}
\hat{\beta} \in [\beta, 1] \quad \text{(perceived present bias)}
\end{equation}

People underestimate their present bias: $\hat{\beta} > \beta$.

\subsubsection{Framework Translation}

Self-knowledge is part of $\Psi_C$:
\begin{equation}
\Psi_T^{effective} = f(\beta, \hat{\beta}, \delta)
\end{equation}

Naïveté ($\hat{\beta} > \beta$) leads to systematic prediction errors.

\subsubsection{Welfare Implications}

With present bias:
\begin{equation}
Q_{experienced} < Q_{long-run optimal}
\end{equation}

People make choices their ``long-run self'' would not endorse.

%==============================================================================
% PART IV: SELF-CONTROL AND COMMITMENT
%==============================================================================
\subsection{Part IV: Self-Control and Commitment Devices}

\subsubsection{Thaler \& Shefrin: Planner-Doer}

Thaler \& Shefrin (1981) modeled internal conflict:

\begin{equation}
\boxed{
\text{Self} = \text{Planner (patient)} + \text{Doer (myopic)}
}
\label{eq:planner-doer}
\end{equation}

\paragraph{The Conflict.}
\begin{itemize}
    \item Planner: Maximizes long-run utility
    \item Doer: Maximizes current period utility
    \item Planner must control or constrain Doer
\end{itemize}

\subsubsection{Commitment Devices}

Sophisticates demand commitment:

\begin{equation}
\text{Commitment} = \text{Restrict future choice set to improve } Q
\end{equation}

\paragraph{Examples.}
\begin{itemize}
    \item Illiquid savings (401k, Christmas clubs)
    \item Deadlines and penalties
    \item Ulysses contracts (addiction treatment)
    \item StickK.com, Beeminder
\end{itemize}

\subsubsection{Ashraf et al.: SEED Accounts}

Ashraf, Karlan \& Yin (2006) tested commitment savings:
\begin{itemize}
    \item Offered commitment savings accounts in Philippines
    \item 28\% take-up despite restrictions
    \item Savings increased 82\% for takers
\end{itemize}

\paragraph{Framework Translation.}
Demand for commitment reveals $\beta < 1$ and sophistication.

\subsubsection{Bryan, Karlan \& Nelson: Commitment Survey}

Bryan et al. (2010) surveyed commitment devices:
\begin{equation}
\text{Commitment demand} = f(\beta, \hat{\beta}, \text{stakes})
\end{equation}

\paragraph{Key Findings.}
\begin{itemize}
    \item Commitment devices can improve welfare
    \item Soft commitments (reminders) also effective
    \item Design matters: flexibility vs. commitment trade-off
\end{itemize}

%==============================================================================
% PART V: APPLICATIONS
%==============================================================================
\subsection{Part V: Applications – Savings, Health, Policy}

\subsubsection{Retirement Savings}

Madrian \& Shea (2001) demonstrated power of defaults:

\begin{equation}
\boxed{
\text{Default enrollment} \Rightarrow \text{Participation: 49\%} \to \text{86\%}
}
\label{eq:default}
\end{equation}

\paragraph{Interpretation.}
Present-biased employees intend to save but procrastinate.
Default enrollment overcomes inertia.

\subsubsection{Chetty et al.: Active vs. Passive Saving}

Chetty et al. (2014) distinguished:
\begin{itemize}
    \item \textbf{Active savers}: Respond to incentives (15\% of population)
    \item \textbf{Passive savers}: Follow defaults (85\% of population)
\end{itemize}

\paragraph{Policy Implication.}
Subsidies work for active savers; defaults work for passive savers.

\subsubsection{Health Behaviors}

Present bias explains:
\begin{itemize}
    \item Overeating (immediate pleasure, delayed health cost)
    \item Under-exercise (immediate cost, delayed benefit)
    \item Smoking, drinking, drugs
    \item Medication non-adherence
\end{itemize}

\subsubsection{Addiction}

Gruber \& Köszegi (2001) modeled addiction with $\beta$-$\delta$:
\begin{equation}
\text{Addiction} = \text{Habit formation} + \text{Present bias}
\end{equation}

Cigarette taxes can be \textit{welfare-improving} even for smokers!

%==============================================================================
% PART VI: INTERGENERATIONAL DISCOUNTING
%==============================================================================
\subsection{Part VI: Intergenerational – Climate and Long Run}

\subsubsection{The Discount Rate Debate}

Climate policy critically depends on discount rate:

\begin{equation}
r = \rho + \eta g
\end{equation}

where $\rho$ = pure time preference, $\eta$ = elasticity of marginal utility, $g$ = growth rate.

\subsubsection{Stern vs. Nordhaus}

\begin{center}
\renewcommand{\arraystretch}{1.3}
\begin{tabular}{|l|c|c|}
\hline
& \textbf{Stern (2007)} & \textbf{Nordhaus (2007)} \\
\hline
$\rho$ (pure time preference) & 0.1\% & 1.5\% \\
Discount rate $r$ & $\sim$1.4\% & $\sim$4.5\% \\
Policy implication & Act now & Wait \\
\hline
\end{tabular}
\end{center}

\paragraph{Framework Translation.}
$\Psi_T^{social}$ (social discount rate) is normative choice with huge policy implications.

\subsubsection{Weitzman: Uncertainty and Declining Rates}

Weitzman (2001) showed:
\begin{equation}
\text{Uncertainty about } r \Rightarrow \text{Declining effective discount rate}
\end{equation}

If future discount rates are uncertain, use lower rates for long horizons.

%==============================================================================
% PART VII: MEASUREMENT
%==============================================================================
\subsection{Part VII: Measuring Time Preferences}

\subsubsection{Methods}

\begin{center}
\renewcommand{\arraystretch}{1.3}
\begin{tabular}{|l|l|l|}
\hline
\textbf{Method} & \textbf{Description} & \textbf{Issues} \\
\hline
Monetary choices & Choose \$X today vs. \$Y later & Arbitrage, liquidity \\
Real effort & Work now vs. later & Domain-specific \\
Consumption & Allocate goods over time & Measurement \\
Revealed preference & Savings, health behaviors & Confounds \\
Convex time budgets & Continuous allocation & Complex \\
\hline
\end{tabular}
\end{center}

\subsubsection{Falk et al.: Global Preferences Survey}

Falk et al. (2018) measured patience globally:
\begin{itemize}
    \item 80,000 people, 76 countries
    \item Patience predicts national income
    \item Substantial cross-country variation in $\Psi_T$
\end{itemize}

\subsubsection{Framework Translation}

$\Psi_T$ varies across:
\begin{itemize}
    \item Individuals ($\beta_i$, $\delta_i$)
    \item Cultures (national patience)
    \item Domains (money vs. health)
    \item Life stages (age effects)
\end{itemize}

%==============================================================================
% PART VIII: SYNTHESIS
%==============================================================================
\subsection{Part VIII: Synthesis – Temporal Context Equations}

\subsubsection{Complete Framework Translation}

\begin{center}
\renewcommand{\arraystretch}{1.5}
\begin{tabular}{|l|c|l|}
\hline
\textbf{Concept} & \textbf{Equation} & \textbf{$\Psi_T$ Element} \\
\hline
Exponential discounting & $U = \sum \delta^t u(c_t)$ & $\Psi_T = \delta$ \\
\hline
Hyperbolic discounting & $D(t) = 1/(1+kt)$ & $\Psi_T(t)$ time-varying \\
\hline
$\beta$-$\delta$ model & $U = u_0 + \beta \sum \delta^t u_t$ & $\Psi_T = (\beta, \delta)$ \\
\hline
Present bias & $\beta < 1$ & Immediate gratification \\
\hline
Sophistication & $\hat{\beta} = \beta$ & Self-knowledge \\
\hline
Naïveté & $\hat{\beta} > \beta$ & Overconfidence \\
\hline
Commitment & Restrict future $C$ & $\Psi_{design}$ for $\Psi_T$ \\
\hline
Social discount rate & $r = \rho + \eta g$ & $\Psi_T^{social}$ for policy \\
\hline
\end{tabular}
\end{center}

\subsubsection{$\Psi_T$ in the Full Framework}

\begin{equation}
\boxed{
C^* = C^*(\Psi_T, \Psi_C, \Psi_I, \ldots) \quad \text{where } \Psi_T = (\beta, \delta, \hat{\beta})
}
\end{equation}

Temporal context interacts with other $\Psi$ dimensions:
\begin{itemize}
    \item $\Psi_T \times \Psi_C$: Cognitive load worsens self-control
    \item $\Psi_T \times \Psi_I$: Institutions can provide commitment
    \item $\Psi_T \times \Psi_E$: Poverty intensifies present bias
\end{itemize}

\subsubsection{Welfare Implications}

With present bias:
\begin{equation}
Q^{long-run} > Q^{chosen} \quad \text{(internality)}
\end{equation}

This justifies:
\begin{itemize}
    \item Nudges and defaults
    \item Commitment devices
    \item Pigouvian taxes on ``sin goods''
    \item Paternalistic policies (if welfare-improving)
\end{itemize}

%==============================================================================
% PART IX: COMPLETE PAPER DOCUMENTATION
%==============================================================================
\subsection{Part IX: Complete Paper Documentation}

%--- FOUNDATIONS ---
\subsubsection{Foundations}

\paragraph{Paper 1: Samuelson (1937)}

``A Note on Measurement of Utility.''
\textit{Review of Economic Studies}, 4(2), 155--161.

\textbf{Summary.}
Introduces discounted utility model. Exponential discounting as standard.
Samuelson himself noted this was \textit{not} meant as accurate description.

\textbf{Framework Translation.}
\begin{equation}
\Psi_T^{standard} = \delta \quad \text{(constant discount factor)}
\end{equation}

\textbf{Key Finding.}
Foundational but problematic as descriptive model.

\paragraph{Paper 2: Strotz (1956)}

``Myopia and Inconsistency in Dynamic Utility Maximization.''
\textit{Review of Economic Studies}, 23(3), 165--180.

\textbf{Summary.}
First formal analysis of time inconsistency. Non-exponential discounting 
leads to changing plans. Proposes commitment as solution.

\textbf{Framework Translation.}
\begin{equation}
\text{Non-exponential } \Psi_T \Rightarrow \text{Time inconsistency}
\end{equation}

\textbf{Key Finding.}
Foundational for behavioral time preferences.

\paragraph{Paper 3: Frederick, Loewenstein \& O'Donoghue (2002)}

``Time Discounting and Time Preference: A Critical Review.''
\textit{Journal of Economic Literature}, 40(2), 351--401.

\textbf{Summary.}
Comprehensive survey of time preference research. Documents anomalies, 
reviews hyperbolic discounting, discusses measurement.

\textbf{Framework Translation.}
\begin{equation}
\Psi_T \neq \delta \quad \text{(single parameter insufficient)}
\end{equation}

\textbf{Key Finding.}
Definitive survey. 6,000+ citations.

%--- HYPERBOLIC DISCOUNTING ---
\subsubsection{Hyperbolic Discounting}

\paragraph{Paper 4: Ainslie (1975)}

``Specious Reward: A Behavioral Theory of Impulsiveness and Impulse Control.''
\textit{Psychological Bulletin}, 82(4), 463--496.

\textbf{Summary.}
Introduces hyperbolic discounting from animal behavior research.
Shows preference reversals empirically.

\textbf{Framework Translation.}
$\Psi_T$ has hyperbolic shape from psychology.

\paragraph{Paper 5: Loewenstein \& Prelec (1992)}

``Anomalies in Intertemporal Choice: Evidence and an Interpretation.''
\textit{Quarterly Journal of Economics}, 107(2), 573--597.

\textbf{Summary.}
Documents anomalies: magnitude effect, sign effect, delay-speedup asymmetry.
Proposes generalized hyperbolic discounting.

\textbf{Framework Translation.}
Multiple anomalies require richer $\Psi_T$.

\paragraph{Paper 6: Laibson (1997)}

``Golden Eggs and Hyperbolic Discounting.''
\textit{Quarterly Journal of Economics}, 112(2), 443--477.

\textbf{Summary.}
Introduces $\beta$-$\delta$ model. Explains illiquid assets as commitment.
Provides tractable framework for applications.

\textbf{Framework Translation.}
\begin{equation}
\Psi_T = (\beta, \delta) \quad \text{with } \beta \approx 0.7
\end{equation}

\textbf{Key Finding.}
4,000+ citations. Standard behavioral discounting model.

%--- SOPHISTICATION AND NAIVETE ---
\subsubsection{Sophistication and Naïveté}

\paragraph{Paper 7: O'Donoghue \& Rabin (1999)}

``Doing It Now or Later.''
\textit{American Economic Review}, 89(1), 103--124.

\textbf{Summary.}
Analyzes sophisticates vs. naïfs. Both procrastinate but for different 
reasons. Sophisticates may preproperate (do things too early).

\textbf{Framework Translation.}
\begin{equation}
\hat{\beta} = \beta \text{ (sophisticated)} \quad \text{vs.} \quad \hat{\beta} = 1 \text{ (naïve)}
\end{equation}

\textbf{Key Finding.}
3,500+ citations. Foundational for behavioral economics.

\paragraph{Paper 8: O'Donoghue \& Rabin (2001)}

``Choice and Procrastination.''
\textit{Quarterly Journal of Economics}, 116(1), 121--160.

\textbf{Summary.}
Introduces partial naïveté. People underestimate but don't ignore 
their present bias. $\hat{\beta} \in (\beta, 1)$.

\textbf{Framework Translation.}
Self-knowledge as continuous parameter.

%--- SELF-CONTROL ---
\subsubsection{Self-Control and Commitment}

\paragraph{Paper 9: Thaler \& Shefrin (1981)}

``An Economic Theory of Self-Control.''
\textit{Journal of Political Economy}, 89(2), 392--406.

\textbf{Summary.}
Planner-doer model of internal conflict. Self-control as agency problem 
within person. Explains mental accounting, rules.

\textbf{Framework Translation.}
\begin{equation}
\text{Self} = \text{Planner}(\delta) + \text{Doer}(\beta)
\end{equation}

\textbf{Key Finding.}
3,000+ citations. Nobel Prize 2017 (Thaler).

\paragraph{Paper 10: Ashraf, Karlan \& Yin (2006)}

``Tying Odysseus to the Mast: Evidence from a Commitment Savings Product.''
\textit{Quarterly Journal of Economics}, 121(2), 635--672.

\textbf{Summary.}
Field experiment on commitment savings in Philippines.
28\% take-up, 82\% savings increase for takers.

\textbf{Framework Translation.}
Revealed preference for commitment = revealed $\beta < 1$.

\textbf{Key Finding.}
1,500+ citations. Classic field experiment.

\paragraph{Paper 11: Bryan, Karlan \& Nelson (2010)}

``Commitment Devices.''
\textit{Annual Review of Economics}, 2, 671--698.

\textbf{Summary.}
Survey of commitment devices in economics. Reviews evidence, 
discusses design principles.

\textbf{Framework Translation.}
$\Psi_{design}$ for self-control.

%--- APPLICATIONS ---
\subsubsection{Applications}

\paragraph{Paper 12: Madrian \& Shea (2001)}

``The Power of Suggestion: Inertia in 401(k) Participation and Savings Behavior.''
\textit{Quarterly Journal of Economics}, 116(4), 1149--1187.

\textbf{Summary.}
Default enrollment in 401(k) dramatically increases participation.
49\% → 86\% participation with automatic enrollment.

\textbf{Framework Translation.}
Defaults overcome present-biased procrastination.

\textbf{Key Finding.}
4,000+ citations. Changed pension policy worldwide.

\paragraph{Paper 13: Chetty et al. (2014)}

``Active vs. Passive Decisions and Crowd-Out in Retirement Savings.''
\textit{Quarterly Journal of Economics}, 129(3), 1141--1219.

\textbf{Summary.}
85\% of savers are passive (follow defaults).
15\% are active (respond to incentives).
Different policies work for each group.

\textbf{Framework Translation.}
Heterogeneous $\Psi_T$ requires heterogeneous policy.

\paragraph{Paper 14: Gruber \& Köszegi (2001)}

``Is Addiction `Rational'? Theory and Evidence.''
\textit{Quarterly Journal of Economics}, 116(4), 1261--1303.

\textbf{Summary.}
Applies $\beta$-$\delta$ to addiction. Cigarette taxes can be welfare-improving 
even from smokers' perspective.

\textbf{Framework Translation.}
\begin{equation}
Q^{long-run}_{smoker} > Q^{chosen} \Rightarrow \text{Tax improves welfare}
\end{equation}

%--- INTERGENERATIONAL ---
\subsubsection{Intergenerational Discounting}

\paragraph{Paper 15: Stern (2007)}

\textit{The Economics of Climate Change: The Stern Review.}
Cambridge University Press.

\textbf{Summary.}
Climate change economics. Uses low discount rate ($\rho \approx 0.1\%$).
Argues for immediate strong action.

\textbf{Framework Translation.}
$\Psi_T^{social} = $ low $\rho \Rightarrow$ aggressive climate policy.

\paragraph{Paper 16: Nordhaus (2007)}

``A Review of the Stern Review on the Economics of Climate Change.''
\textit{Journal of Economic Literature}, 45(3), 686--702.

\textbf{Summary.}
Critique of Stern. Argues for higher discount rate based on market returns.
Different $\Psi_T^{social}$ leads to different policy.

\textbf{Framework Translation.}
$\Psi_T^{social}$ is fundamental normative choice.

\textbf{Key Finding.}
Nobel Prize 2018 (Nordhaus).

\paragraph{Paper 17: Weitzman (2001)}

``Gamma Discounting.''
\textit{American Economic Review}, 91(1), 260--271.

\textbf{Summary.}
Uncertainty about discount rates implies declining discount rates.
Long-term projects should use lower rates.

\textbf{Framework Translation.}
$\Psi_T$ uncertainty → use lower rates for long horizons.

%--- MEASUREMENT ---
\subsubsection{Measurement and Heterogeneity}

\paragraph{Paper 18: Falk et al. (2018)}

``Global Evidence on Economic Preferences.''
\textit{Quarterly Journal of Economics}, 133(4), 1645--1692.

\textbf{Summary.}
Measures patience and other preferences in 76 countries.
Patience predicts income, patience varies substantially.

\textbf{Framework Translation.}
$\Psi_T$ varies systematically across countries.

\textbf{Key Finding.}
GPS dataset used widely.

\paragraph{Paper 19: Andreoni \& Sprenger (2012)}

``Estimating Time Preferences from Convex Budgets.''
\textit{American Economic Review}, 102(7), 3333--3356.

\textbf{Summary.}
New method using convex time budgets. Separates discounting from 
utility curvature. Estimates $\beta \approx 1$ (challenges present bias).

\textbf{Framework Translation.}
$\Psi_T$ measurement methodology.

\paragraph{Paper 20: Cohen et al. (2020)}

``Measuring Time Preferences.''
\textit{Journal of Economic Literature}, 58(2), 299--347.

\textbf{Summary.}
Comprehensive survey of time preference measurement.
Reviews methods, findings, and controversies.

\textbf{Framework Translation.}
State of the art on $\Psi_T$ measurement.

%%%%%%%%%%%%%%%%%%%%%%%%%%%%%%%%%%%%%%%%%%%%%%%%%%%%%%%%%%%%%%%%%%%%%%%%%%%%%%%
\subsection{Citation and Verification Note}

\textbf{Nobel Prizes represented:}
\begin{itemize}
    \item Samuelson (1970) ``for scientific work developing static and dynamic economic theory''
    \item Thaler (2017) ``for contributions to behavioral economics''
    \item Nordhaus (2018) ``for integrating climate change into long-run macroeconomic analysis''
\end{itemize}

Combined Google Scholar citations for the 20 papers: $>$50,000.

This appendix provides the theoretical and empirical foundations for the
temporal dimension $\Psi_T$ of the framework, showing how time preferences
fundamentally shape economic behavior and policy.


%%%%%%%%%%%%%%%%%%%%%%%%%%%%%%%%%%%%%%%%%%%%%%%%%%%%%%%%%%%%%%%%%%%%%%%%%%%%%%%
\section{Key Results}
\label{sec:ctm-results}
%%%%%%%%%%%%%%%%%%%%%%%%%%%%%%%%%%%%%%%%%%%%%%%%%%%%%%%%%%%%%%%%%%%%%%%%%%%%%%%

The 20 papers reviewed establish:

\begin{enumerate}
    \item \textbf{Hyperbolic discounting} dominates exponential discounting empirically ($\beta \approx 0.7$)
    \item \textbf{Naivete} about time inconsistency is common and welfare-reducing
    \item \textbf{Commitment devices} can improve outcomes for sophisticated agents
    \item \textbf{Cross-country variation} in patience is substantial and predictive
    \item \textbf{Measurement methodology} continues to evolve (convex budgets, CTB)
\end{enumerate}


%%%%%%%%%%%%%%%%%%%%%%%%%%%%%%%%%%%%%%%%%%%%%%%%%%%%%%%%%%%%%%%%%%%%%%%%%%%%%%%
\section{Summary}
\label{sec:ctm-summary}
%%%%%%%%%%%%%%%%%%%%%%%%%%%%%%%%%%%%%%%%%%%%%%%%%%%%%%%%%%%%%%%%%%%%%%%%%%%%%%%

\begin{tcolorbox}[colback=gray!5!white, colframe=gray!75!black,
    title=Appendix UNMAPPED_CTM Summary: Temporal Context]

\textbf{Core Insight:}
Time preferences ($\Psi_T$) are a fundamental dimension of context that shapes economic behavior. Hyperbolic discounting and present bias are robust empirical phenomena with important policy implications.

\textbf{Key Papers:}
\begin{itemize}[nosep]
    \item Samuelson (1937): Foundation of discounted utility
    \item Laibson (1997): $\beta\delta$ model of hyperbolic discounting
    \item O'Donoghue \& Rabin (1999, 2001): Sophistication vs. naivete
    \item Nordhaus (1994): Intergenerational discounting for climate
\end{itemize}

\textbf{Framework Integration:}
$\Psi_T$ integrates with the 10C CORE through the WHEN dimension, providing theoretical and empirical foundations for temporal context effects.

\end{tcolorbox}


%%%%%%%%%%%%%%%%%%%%%%%%%%%%%%%%%%%%%%%%%%%%%%%%%%%%%%%%%%%%%%%%%%%%%%%%%%%%%%%
\section{References}
\label{sec:ctm-references}
%%%%%%%%%%%%%%%%%%%%%%%%%%%%%%%%%%%%%%%%%%%%%%%%%%%%%%%%%%%%%%%%%%%%%%%%%%%%%%%

\nocite{bcm_master}
See master bibliography (\texttt{bcm\_master.bib}) for complete citations.
For comprehensive symbol definitions, see Appendix G (Master Glossary).

%%%%%%%%%%%%%%%%%%%%%%%%%%%%%%%%%%%%%%%%%%%%%%%%%%%%%%%%%%%%%%%%%%%%%%%%%%%%%%%
